\documentclass[12pt,a4paper]{article}

% ==== PACKAGES ====
\usepackage[utf8]{inputenc}
\usepackage[T1]{fontenc}
\usepackage[english]{babel}
\usepackage{graphicx}
\usepackage{setspace}
\usepackage{geometry}
\geometry{margin=1in}
\usepackage{titling}
\usepackage{lmodern}
\usepackage{float}
\usepackage{caption}
\usepackage{listings}
\usepackage{xcolor}
\usepackage{url}

\lstdefinestyle{dockerstyle}{
	basicstyle=\footnotesize\ttfamily,
	backgroundcolor=\color{gray!10},
	breaklines=true,
	frame=single,
	columns=fullflexible
}

\setlength{\parindent}{0pt}
% ==== DOCUMENT ====
\begin{document}

% ==== TITLE PAGE ====
\begin{titlepage}
    \centering
    \vspace*{2cm}
    
    {\Huge \textbf{Workshop No. 4 - Containerization, Acceptance, and CI/CD}}\\[1.2cm]
    {\LARGE \textbf{Software Engineering Seminar}}\\[0.5cm]
    {\Large \textbf{Universidad Distrital Francisco José de Caldas}}\\[0.2cm]
    {\large \textbf{Systems Engineering Program}}\\[2cm]

    % Optional logo (remove comment if you have a logo)
    % \includegraphics[width=0.25\textwidth]{logo_ud.png}\\[1.5cm]

    {\Large \textbf{Presented by:}}\\[0.5cm]
    {\large Brayan Alejandro Riveros Rodríguez \\ \textit{20201020084}}\\[0.3cm]
    {\large Carlos Andrés Pescador Castro \\ \textit{20182020139}}\\[0.3cm]
    {\large Cristian Santiago Ríos Vásquez \\ \textit{20201020107}}\\[2cm]

    {\Large Bogotá D.C., Colombia}\\[0.3cm]
    {\large November 29, 2025}
    
    \vfill
\end{titlepage}

% ==== CONTENT ====

\section{Containerization}

\subsection{Dockerfiles}
For this project, each backend includes a custom Dockerfile. The goal of these Dockerfiles is to produce lightweight and efficient container images ready for deployment in local and production environments.\\ 

The Java backend Dockerfile uses a multi-stage build. In the first stage, Maven downloads all dependencies and creates the application JAR file. In the second stage, a smaller Amazon Corretto runtime image is used to run the final application. This approach reduces the final image size and improves performance. The container exposes port 8080 and runs the service using a simple java -jar command.\\

\begin{lstlisting}[style=dockerstyle,caption={Dockerfile for Java Backend}]
	FROM maven:3.9.11-amazoncorretto-25-alpine AS build
	WORKDIR /app
	COPY pom.xml .
	RUN mvn dependency:go-offline -B
	COPY src ./src
	RUN mvn clean package -DskipTests
	
	FROM amazoncorretto:25-alpine AS runtime
	WORKDIR /app
	COPY --from=build /app/target/*.jar app.jar
	EXPOSE 8080
	ENTRYPOINT ["java", "-jar", "app.jar"]
\end{lstlisting}
\vspace{0.5cm} 

The Python backend Dockerfile also uses a multi-stage build. The first stage installs dependencies using Poetry and creates a virtual environment inside the project folder. The second stage copies the virtual environment and application code into a clean Python 3.11 image. A non-root user is created for security, and the service is started using Uvicorn on port 8000. A health check is also included to verify that the API is running correctly.\\

\begin{lstlisting}[style=dockerstyle,caption={Dockerfile for Python Backend}]
	FROM python:3.11-slim AS builder
	
	ENV POETRY_VERSION=1.8.3 \
	POETRY_HOME=/opt/poetry \
	POETRY_NO_INTERACTION=1 \
	POETRY_VIRTUALENVS_IN_PROJECT=1 \
	POETRY_VIRTUALENVS_CREATE=1 \
	POETRY_CACHE_DIR=/tmp/poetry_cache
	
	RUN apt-get update && apt-get install -y --no-install-recommends \
	curl gcc libpq-dev && rm -rf /var/lib/apt/lists/*
	
	RUN curl -sSL https://install.python-poetry.org | python3 -
	WORKDIR /app
	
	COPY pyproject.toml poetry.lock* ./
	RUN poetry lock --no-update 2>/dev/null || true && \
	poetry install --only main --no-root --no-directory && \
	rm -rf $POETRY_CACHE_DIR
	
	FROM python:3.11-slim
	ENV PYTHONUNBUFFERED=1 PYTHONDONTWRITEBYTECODE=1 \
	PATH="/app/.venv/bin:$PATH"
	
	RUN apt-get update && apt-get install -y --no-install-recommends \
	libpq5 && rm -rf /var/lib/apt/lists/*
	
	WORKDIR /app
	COPY --from=builder /app/.venv /app/.venv
	COPY . .
	
	RUN useradd -m -u 1000 appuser && \
	chown -R appuser:appuser /app
	USER appuser
	
	EXPOSE 8000
	
	HEALTHCHECK --interval=30s --timeout=10s --start-period=40s --retries=3 \
	CMD python -c "import urllib.request; urllib.request.urlopen('http://localhost:8000/api/v1/health')" || exit 1
	
	CMD ["uvicorn", "src.main:app", "--host", "0.0.0.0", "--port", "8000"]
\end{lstlisting}
\vspace{0.5cm} 

These Dockerfiles allow both services to run consistently across different environments and ensure that all dependencies and runtime configurations are properly isolated inside each container.

\subsection{Docker Compose Configuration}
To orchestrate the application components, we created multiple docker-compose.yml and docker-stack.yml files for development and deployment. These files define how each container interacts with its database and with the reverse proxy when needed.\\ 

For local development, Docker Compose builds the images directly from the Dockerfiles and maps the service ports to the host machine.\\ 

For the Java backend, the compose file includes a MySQL database initialized with SQL scripts.\\ 

\begin{lstlisting}[style=dockerstyle,caption={docker-compose.yml for Java Backend (Local)}]
	version: '3.8'
	services:
	java-backend:
	image: java-backend:latest
	build:
	context: .
	dockerfile: Dockerfile
	container_name: java-backend
	ports:
	- "8080:8080"
	env_file:
	- .env
	depends_on:
	- java-backend-db
	networks:
	- java-backend-network
	
	java-backend-db:
	image: mysql:8.0
	container_name: java-backend-db
	env_file:
	- .env
	volumes:
	- ./init_db_scripts/schema_db.sql:/docker-entrypoint-initdb.d/schema_db.sql
	- db_data:/var/lib/mysql
	healthcheck:
	test: ["CMD", "mysqladmin", "ping", "-h", "localhost", "-u", "${MYSQL_USER}", "-p${MYSQL_PASSWORD}"]
	interval: 10s
	timeout: 5s
	retries: 5
	networks:
	- java-backend-network
	
	volumes:
	db_data:
	
	networks:
	java-backend-network:
	driver: bridge
\end{lstlisting}
\vspace{0.5cm} 

For the Python backend, the compose file includes a PostgreSQL database with its initialization scripts and health checks.\\

\begin{lstlisting}[style=dockerstyle,caption={docker-compose.yml for Python Backend (Local Development)}]
	version: '3.8'
	
	services:
	python-backend-db:
	image: postgres:15-alpine
	container_name: python-backend-db
	restart: unless-stopped
	env_file:
	- .env
	volumes:
	- ./init_db_scripts/schema_db.sql:/docker-entrypoint-initdb.d/schema_db.sql:ro
	- postgres_data:/var/lib/postgresql/data
	healthcheck:
	test: ["CMD-SHELL", "pg_isready -U ${POSTGRES_USER:-postgres} -d ${POSTGRES_DB:-postgres}"]
	interval: 10s
	timeout: 5s
	retries: 5
	
	python-backend:
	build:
	context: .
	dockerfile: Dockerfile
	image: python-backend:latest
	container_name: python-backend
	restart: unless-stopped
	ports:
	- "${BACKEND_PORT:-8000}:8000"
	env_file:
	- .env
	depends_on:
	python-backend-db:
	condition: service_healthy
	command: uvicorn src.main:app --host 0.0.0.0 --port 8000 --reload
	
	volumes:
	postgres_data:
	driver: local
\end{lstlisting}
\vspace{0.5cm} 

Both services load environment variables from .env files and wait for the databases to become healthy before starting.\\

For production, the docker-stack.yml files define the same services but include additional configuration such as replicas, restart policies, placement constraints, and Traefik labels. These labels allow each backend to be published through HTTPS using different hostnames and connected to a shared reverse proxy network.\\ 

\begin{lstlisting}[style=dockerstyle,caption={docker-stack.yml for Java Backend}]
	version: '3.8'
	services:
	java-backend:
	image: java-backend:latest
	env_file:
	- .env
	depends_on:
	- java-backend-db
	networks:
	- java-backend-network
	- reverse_proxy
	deploy:
	replicas: 1
	restart_policy:
	condition: on-failure
	delay: 10s
	max_attempts: 5
	placement:
	constraints:
	- node.role == manager
	labels:
	- "traefik.enable=true"
	- "traefik.docker.network=reverse_proxy"
	- "traefik.http.routers.java-backend.entrypoints=websecure"
	- "traefik.http.routers.java-backend.tls.certresolver=production"
	- "traefik.http.routers.java-backend.rule=Host(`authbackbs.glud.org`)"
	- "traefik.http.services.java-backend.loadbalancer.server.port=8080"
\end{lstlisting}

\vspace{0.5cm}

\begin{lstlisting}[style=dockerstyle,caption={docker-stack.yml for Python Backend (Docker Swarm)}]
	version: '3.8'
	
	services:
	python-backend:
	image: python-backend:latest
	env_file:
	- .env
	depends_on:
	- python-backend-db
	command: uvicorn src.main:app --host 0.0.0.0 --port 8000
	networks:
	- backend-network
	- reverse_proxy
	deploy:
	replicas: 1
	restart_policy:
	condition: on-failure
	placement:
	constraints:
	- node.role == manager
	labels:
	- "traefik.enable=true"
	- "traefik.docker.network=reverse_proxy"
	- "traefik.http.routers.python-backend.entrypoints=websecure"
	- "traefik.http.routers.python-backend.tls.certresolver=production"
	- "traefik.http.routers.python-backend.tls=true"
	- "traefik.http.routers.python-backend.rule=Host(`coursefinderbs.glud.org`)"
	- "traefik.http.services.python-backend.loadbalancer.server.port=8000"
	
	python-backend-db:
	image: postgres:15-alpine
	env_file:
	- .env
	networks:
	- backend-network
	volumes:
	- ./init_db_scripts/schema_db.sql:/docker-entrypoint-initdb.d/schema_db.sql:ro
	- postgres_data:/var/lib/postgresql/data
	healthcheck:
	test: ["CMD-SHELL", "pg_isready -U ${POSTGRES_USER:-postgres}"]
	interval: 10s
	timeout: 5s
	retries: 5
	deploy:
	replicas: 1
	restart_policy:
	condition: on-failure
	placement:
	constraints:
	- node.role == manager
	
	volumes:
	postgres_data:
	driver: local
	
	networks:
	backend-network:
	driver: overlay
	reverse_proxy:
	external: true
\end{lstlisting}
\vspace{0.5cm} 

Using Docker Compose and Docker Swarm makes it possible to run the system as a group of independent services while keeping the configuration simple, repeatable, and easy to deploy.

\subsection{Makefile Usage}
The backends also includes a Makefile to simplify common tasks. This file includes commands for installing dependencies, running tests, linting the code, cleaning temporary files, building Docker images, and deploying the stack in Docker Swarm. Those Makefile helps maintain a clean and consistent workflow for all team members.\\

\begin{lstlisting}[style=dockerstyle,caption={Makefile for Java Backend}]	
	install: 
	mvn clean package
	
	dev:
	mvn clean spring-boot:run
	
	build: 
	docker build -t java-backend:latest .
	
	build-no-cache: 
	docker build --no-cache -t java-backend:latest .
	
	compose-up:  
	docker compose up -d
	
	compose-up-build:  
	docker compose up -d --build
	
	compose-down: 
	docker compose down
	
	stack-deploy: 
	docker stack deploy -c docker-stack.yml java-backend
	
	stack-remove: 
	docker stack rm java-backend
	
	stack-ps: 
	docker stack ps java-backend
	
	stack-services:
	docker stack services java-backend
\end{lstlisting}

\vspace{0.5cm}

\begin{lstlisting}[style=dockerstyle,caption={Makefile for Python Backend}]
	install:
	poetry install
	
	install-prod:
	poetry install --only main
	
	update:
	poetry update
	
	lock:
	poetry lock
	
	dev:
	poetry run uvicorn src.main:app --reload --host 0.0.0.0 --port 8000
	
	test:
	poetry run pytest
	
	test-cov:
	poetry run pytest --cov=src --cov-report=html --cov-report=term
	
	lint:
	poetry run ruff check src
	
	lint-fix:
	poetry run ruff check --fix src
	
	format:
	poetry run ruff format src
	
	type-check:
	poetry run mypy src
	
	clean:
	find . -type d -name "__pycache__" -exec rm -rf {} +
	rm -rf htmlcov .coverage
	
	build:
	docker build -t python-backend:latest .
	
	build-no-cache:
	docker build --no-cache -t python-backend:latest .
	
	compose-up:
	docker compose up -d
	
	compose-up-build:
	docker compose up -d --build
	
	compose-down:
	docker compose down
	
	stack-deploy:
	docker stack deploy -c docker-stack.yml python-backend
	
	stack-remove:
	docker stack rm python-backend
	
	stack-ps:
	docker stack ps python-backend
	
	stack-services: 
	docker stack services python-backend
	
\end{lstlisting}

\subsection{Summary}

In summary, the containerization process provides a complete and isolated environment for both backends. 
The Dockerfiles create optimized and lightweight images, while the Compose and Stack files configure the services, databases, networks, and the reverse proxy for different environments.\\ 

In addition to this, both backends include Makefiles that automate common tasks such as installing dependencies, running the application in development mode, building Docker images, and deploying the services with Docker Swarm. 
These Makefiles help maintain a consistent workflow and make the development and deployment process faster and easier. 
Overall, these containers make the project simple to run, test, and deploy in both development and production environments.

\section{Acceptance Tests with Cucumber}

\subsection{Overview}
For this project, we used Cucumber to create acceptance tests based on the main user stories. The goal of these tests is to verify the system behavior from the point of view of the final user. Cucumber allows us to write tests in a natural language format (Gherkin), making the scenarios easy to understand and simple to maintain.

The acceptance tests help validate that the backends return the correct responses and that the main flows of the system work as expected. Each feature file represents a specific user story, and each scenario describes one possible interaction or outcome.

\subsection{Feature Files}

The feature files were written using the Gherkin syntax. Each file includes \textit{Given}, \textit{When}, and \textit{Then} steps that describe the expected behavior. The scenarios focus on validating actions such as searching data, creating records, authenticating users, or interacting with the REST API.\\ 

Below is a general description of the structure used in the feature files:

\begin{itemize}
	\item \textbf{Given:} Defines the initial state or preconditions. For example, the user is authenticated or the database contains specific records.
	\item \textbf{When:} Describes the action performed by the user, such as sending a request to an API endpoint.
	\item \textbf{Then:} Defines the expected result, such as receiving a correct HTTP status or a JSON response with the expected fields.
\end{itemize}

This structure helps ensure clarity and consistency across all acceptance tests.

\subsection{Step Definitions}

For each feature file, we implemented step definitions that connect the Gherkin steps with executable test code. These step definitions send HTTP requests to the backend, receive the responses, and verify the expected results.\\ 

The step definitions were organized by functionality. For example:

\begin{itemize}
	\item Steps for API requests and responses.
	\item Steps for authentication or user validation.
	\item Steps for comparing JSON fields or checking HTTP status codes.
\end{itemize}

This modular structure makes the test suite easier to extend and maintain.

\subsection{Execution and Test Results}

The Cucumber tests were executed using the project's testing environment. After running the test suite, Cucumber generated a report showing which scenarios passed or failed.\\ 

The results showed the following:

\begin{itemize}
	\item The main user stories behaved correctly.
	\item The REST API returned the expected responses for valid scenarios.
	\item Error scenarios were handled correctly, such as missing parameters or invalid requests.
\end{itemize}

The reports also include execution time and a breakdown of each scenario. These results help confirm that both backends meet the acceptance criteria of the project.

\subsection{Summary}

The acceptance tests created with Cucumber verify the correct behavior of the application from the user perspective. The feature files, step definitions, and reports provide a clear and structured validation of the system. This ensures that the project satisfies the main user requirements and that the core functionality works as expected.



\section{JMeter Test Plans and Results}

\subsection{Overview}

For this project, we used Apache JMeter to perform stress testing on the REST APIs of both backends. The goal of these tests was to evaluate how the system behaves under different levels of load and to identify possible performance problems. JMeter allows us to simulate many users sending requests at the same time, which helps us understand the stability and responsiveness of the application.

\subsection{Test Plan Design}

The JMeter test plans were created to cover the most important API endpoints. Each plan includes several elements such as:

\begin{itemize}
	\item \textbf{Thread Groups:} Used to simulate multiple virtual users. We configured different numbers of users to test the system under light, medium, and heavy load.
	\item \textbf{HTTP Samplers:} These samplers send requests to the backend endpoints. Each sampler represents a specific API action such as creating, reading, or searching data.
	\item \textbf{Assertions:} Added to verify that each response contains the correct status code and expected fields.
	\item \textbf{Timers:} Used to simulate realistic delays between requests.
\end{itemize}

This structure ensures that the tests follow an organized and repeatable pattern.

\subsection{Scenarios Executed}

The following types of scenarios were tested:

\begin{itemize}
	\item \textbf{Single Endpoint Load:} Stress tests for individual endpoints to check their response time and stability.
	\item \textbf{Multiple Endpoint Sequence:} Simulates a user flow with several consecutive operations.
	\item \textbf{High Concurrency Scenario:} Simulates many users working at the same time to evaluate how the system behaves under heavy load.
\end{itemize}

These scenarios help verify both isolated and combined behavior of the APIs.

\subsection{Results and Observations}

After executing the test plans, JMeter generated detailed reports. The results showed the following:

\begin{itemize}
	\item The system maintained stable response times under normal and moderate load.
	\item No critical errors appeared during concurrent requests.
	\item Response times increased under heavy load, which is expected in stress testing.
	\item The backends correctly handled all tested endpoints.
\end{itemize}

We also reviewed metrics such as response time distribution, throughput, and error rate. These results help understand the system capacity and potential areas for optimization.

\subsection{Summary}

The JMeter stress tests confirmed that the application can handle different levels of load while keeping stable behavior. The test plans provide a clear view of how the APIs respond in realistic scenarios, and the results help validate that the system is ready for production-like environments.

\section{GitHub Actions Workflow}

\subsection{Overview}

For this project, we created a GitHub Actions workflow to automate important tasks such as running tests and building Docker images. The goal of this workflow is to support continuous integration and continuous delivery (CI/CD). Each time we push new code to the repository, the workflow runs automatically and checks that the system continues to work correctly.

\subsection{Workflow Structure}

The workflow file includes several main steps:

\begin{itemize}
	\item \textbf{Code Checkout:} Downloads the project from the repository so the workflow can work with the files.
	\item \textbf{Environment Setup:} Installs the correct version of Java, Python, or other tools required by the backends.
	\item \textbf{Dependency Installation:} Installs all project dependencies so the tests can run correctly.
	\item \textbf{Test Execution:} Runs the unit tests or acceptance tests defined in the project.
	\item \textbf{Docker Build:} Builds the Docker images for the Java backend, the Python backend, or both.
\end{itemize}

This structure ensures that each push or pull request is automatically validated.

\subsection{Automatic Testing}

One important part of the workflow is the automatic test execution. When the workflow runs, it executes the corresponding test suites. If any test fails, the workflow stops and reports the error. This helps identify problems early and ensures code quality.

The results show whether all tests passed and allow developers to correct errors before merging new changes.

\subsection{Docker Image Build}

The workflow also builds the Docker images for the backends. This guarantees that the Dockerfiles are correct and that the application can be containerized without errors. The images can later be published or used in a deployment environment.

Building the images automatically helps maintain consistency between development and production environments.

\subsection{Workflow Results}

GitHub Actions provides logs and status messages for each step of the workflow. These messages indicate:

\begin{itemize}
	\item if the tests passed correctly,
	\item if the Docker images were built without errors,
	\item and if the full workflow completed successfully.
\end{itemize}

These results confirm that the system is stable and that the CI/CD process is working as expected.

\subsection{Summary}

The GitHub Actions workflow gives the project an automated process for testing and building the application. This CI/CD setup improves code quality, reduces manual work, and ensures that every update to the repository is validated. It also enables faster and more reliable deployments for both backends.


\section{Conclusion}


\end{document}